\documentclass{jarticle}
\usepackage[dvipdfmx]{graphicx}
\topmargin       -10.0mm
\oddsidemargin   -0.4mm
\evensidemargin  -0.4mm
\headheight       0.0mm
\headsep          0.0mm
\footskip        10.0mm
%
\textheight     260.0mm
\textwidth      170.0mm
%
\setlength{\parindent}{1zw}
%
\usepackage{graphicx}
\usepackage{multirow}
\usepackage{wrapfig}
\makeatletter
\def \fgcaption{\def \@captype{figure} \caption}
\def \tbcaption{\def \@captype{table} \caption}
\makeatother
\def \figurename{Fig.}
\def \tablename{Table}
%
\def\topfraction{.99}
\def\bottomfraction{.99}
\def\textfraction{.01}
\def\floatpagefraction{.99}
%
\pagestyle{empty}
%
\begin{document}
%
\begin{center}
 \begin{tabular}{|c|c|c|} \hline
  \parbox{50mm}{\vspace{1mm} \centering{\Large\bf 2026/06/11}} & 
  \parbox{55mm}{\vspace{1mm} \centering{\Large\bf B4　堀田裕樹}} & 
  \parbox{50mm}{\vspace{1mm} \centering{\Large\bf No.~9}}
  \\ [2mm] \hline
%
  \multicolumn{3}{|c|}{%
   \parbox{155mm}{\centering\Large\bf
    \vspace{1mm}
    widthの細分化シミュレーションによる連続的な特性変化の評価
   }
  }
  \\ [2mm] \hline\hline
%
  \multicolumn{3}{|c|}{%
   \parbox{155mm}{\large\vspace{1mm}
%
    {\Large\bf 今月の目標}
%

%
    {\Large\bf 進捗の要点}
%
    \begin{itemize}
     \item widthのシミュレーション間隔を0.05mm刻みに細分化（0.15mm--1.00mm）
     \item 波形およびFFTスペクトルの変化をアニメーション(GIF)化し，連続的な遷移を確認
    \end{itemize}
%
    {\Large\bf 今後の実施事項}
%
    \begin{itemize}
     \item 実験データとの詳細なマッチング
    \end{itemize}
%
   }
  }
  \\ [2mm] \hline
 \end{tabular}
%
\end{center}
%
\par
先週のシミュレーションでは，widthの変化間隔が大きく，特性変化の連続性を把握することが困難であった．そこで本週は，widthを 0.15mm から 1.00mm まで 0.05mm 刻みでシミュレーションを実施した．

動的な変化を視覚的に評価するため，波形およびFFTの遷移をGIFアニメーションとして作成した．

これにより,widthの増加につれて波形の振幅が徐々に減容する様子が確認できた.またテイル波の先頭部分に関しては増減が確認された.
スペクトルに関しても,widthの増加につれて振幅が減容する様子が確認できた.またディップの位置が移動したり消失が確認された.
\begin{center}
  \begin{minipage}[b]{0.49\linewidth}
    \centering
    \includegraphics[width=6.8cm]{../../data_im/comparison_plots/comparison_waveform.png}
    \vspace{-2mm}
    \fgcaption{Waveform (Representative: Pitch 2000)}
  \end{minipage}
  \begin{minipage}[b]{0.49\linewidth}
    \centering
    \includegraphics[width=6.8cm]{../../data_im/comparison_plots/comparison_fft.png}
    \vspace{-2mm}
    \fgcaption{FFT (Representative: Pitch 2000)}
  \end{minipage}
  \label{fig:results}
\end{center}

%
\end{document}
